\chapter{Infinitesimal deformations}
\label{ch:deformations}

We first work with square-zero extensions of rings and chosen
identifications of their kernels. We then specialize to local Artin
\(k\)-algebras with a specified residue field \(k\).
The construction extends the linear lifting equation
\cite[FL-06]{core} from lifting a map into a given ring to
constructing the target extension itself.

\section{Lifting a ring over a prescribed extension}

\begin{theorem}[Square-zero extensions from a presentation]\label{res:T01}
Let \(A'\twoheadrightarrow A\) have square-zero kernel \(I\).
Let \(B\) be an \(A\)-algebra, \(N\) a \(B\)-module, and \(c:I\to N\)
an \(A\)-linear map. A lift is an \(A'\)-algebra \(B'\twoheadrightarrow B\)
with square-zero kernel identified with \(N\), inducing \(c\) on \(I\).
Take a polynomial presentation \(P'=A'[X_i]\twoheadrightarrow B\)
with kernel \(Q\). Let \(\overline P=P'/IP'\), let
\(J=\ker(\overline P\to B)\), and write \(K\) for the image of
\(IP'\to Q/Q^2\).
A lift exists if and only if the induced \(P'\)-linear map
\(IP'\to N\) factors through \(K\) and the resulting map
\(K\to N\) extends to \(Q/Q^2\). Once the factorization condition
holds, the extension condition is the vanishing of its connecting
class in \(\operatorname{Ext}^1_B(J/J^2,N)\).
\end{theorem}
\begin{proof}
Given a lift, choose lifts of the variable images in \(B'\).
The resulting map \(P'\to B'\) sends \(Q\) into \(N\), kills
\(Q^2\), and induces a \(B\)-linear homomorphism \(\rho:Q/Q^2\to N\)
extending the prescribed map. Conversely, make
\(N\) a module over \(P'/Q^2\) through \(B\) and form the trivial
square-zero extension \((P'/Q^2)\ltimes N\), with multiplication
\((a,u)(b,v)=(ab,av+bu)\).
The graph \(\{(q,-\rho(q)):q\in Q/Q^2\}\) is an ideal because
\(Q/Q^2\) acts trivially on \(N\). Quotienting identifies the kernel
of the map to \(B\) with \(N\), via \((q,n)\mapsto\rho(q)+n\).
The induced map on \(I\) is \(c\). This construction is inverse
to the previous one with lifted variables because the map to \(B'\)
is surjective and has exactly that graph as its kernel.

The quotient \((Q/Q^2)/K\) identifies with \(J/J^2\), since \(Q\)
surjects to \(J\) with kernel \(IP'\) and its square has image \(J^2\).
Apply \(\operatorname{Hom}_B(-,N)\) to
\(0\to K\to Q/Q^2\to J/J^2\to0\).
Exactness identifies extension of the prescribed map with zero
connecting class. Thus the vanishing test is exact for each
presentation. Its truth is independent of the presentation since
it is equivalent to the existence of a lift.
\end{proof}

\begin{proposition}[Differences and infinitesimal automorphisms]\label{res:T02}
If the lifts in Result~\ref{res:T01} exist, their isomorphism
classes form a torsor under
\[
 T^1(B/A;N)=
 \operatorname{coker}\bigl(
  \operatorname{Hom}_B(\Omega_{\overline P/A}\otimes_{\overline P}B,N)
  \longrightarrow \operatorname{Hom}_B(J/J^2,N)\bigr).
\]
Here isomorphisms preserve \(B,N\) and the \(A'\)-structure.
The automorphisms of a chosen lift inducing the identity on
\(B,N\) are naturally \(\operatorname{Der}_A(B,N)\).
The displayed group is
\(\operatorname{Ext}^1_B([J/J^2\to\Omega_{\overline P/A}\otimes B],N)\)
where the complex is in degrees \(-1,0\).
\end{proposition}
\begin{proof}
Any two extensions of the prescribed map \(\rho|_K\) differ by a
map \(J/J^2\to N\). Changing the chosen variable lifts by elements
of \(N\) changes \(\rho\) by the restriction of a polynomial
derivation. That restriction factors through the conormal map.
An isomorphism of two lifts sends one list of variable lifts to
another such list, so this equivalence accounts for precisely the
isomorphic lifts, by the explicit graph construction.
An automorphism with the specified identities has the form
\(z\mapsto z+d(\bar z)\) with \(d:B\to N\). Additivity and
multiplicativity are equivalent to additivity and the Leibniz
rule for \(d\). Respecting the \(A'\)-map makes \(d\) an
\(A\)-derivation. Conversely this formula always gives the required
automorphism with inverse obtained by replacing \(d\) by \(-d\).

For the last formula, use the triangle with degree-zero free
term \(\Omega_{\overline P/A}\otimes B\), the two-term complex,
and the quotient complex \(J/J^2[1]\). Apply the derived Hom
long exact sequence in \cite[R11]{homological}. The degree-one
Ext from the free term is zero, and the degree-one Ext from
\(J/J^2[1]\) is \(\operatorname{Hom}_B(J/J^2,N)\).
Thus degree one is exactly the indicated cokernel.
\end{proof}

\begin{lemma}[The flat-lift condition over Artin rings]\label{res:T03}
Let \(A'\to A\) be a square-zero surjection of local Artin
\(k\)-algebras with identity on \(k\). Let \(B_A\) be a flat
\(A\)-algebra. A lift of \(B_A\) which is \(A'\)-flat and reduces
to \(B_A\) gives a solution of Result~\ref{res:T01} with
\(N=I\otimes_A B_A\) and \(c(i)=i\otimes1\). Conversely a
solution for this \(N,c\) is flat and reduces to \(B_A\).
\end{lemma}
\begin{proof}
A flat module \(M\) over a local Artin ring is free. Indeed lifts
of a basis modulo the maximal ideal give a surjection from a free
module by nilpotence. Flatness implies its kernel is zero modulo
the maximal ideal after tensoring the sequence with the residue
field. Nilpotence gives zero kernel even for infinite bases.
For a flat lift, tensor \(0\to I\to A'\to A\to0\) with the lifted
algebra to obtain exactly the stated sequence with
\(I\otimes_{A'} B'=I\otimes_A B_A\).

Conversely let \(B'\) be a solution. Lift an \(A\)-basis of \(B_A\)
to \(B'\) and compare the exact sequence with
\(0\to I^{(S)}\to (A')^{(S)}\to A^{(S)}\to0\), indexed by the basis.
The map on left terms is the canonical basis isomorphism since
the map on \(I\) is \(c\), and the map on right terms is an
isomorphism by choice. Injectivity and surjectivity of the middle
map now follow by chasing the sequence. In particular \(B'\)
is free and its base change is the prescribed \(B_A\).
\end{proof}

\section{Complete intersections and an obstruction example}

\begin{theorem}[Unobstructed polynomial complete intersections]\label{res:T04}
Let \(P=k[x_1,\ldots,x_n]\), \(f_1,\ldots,f_s\) be a regular
sequence, and \(B=P/(f_1,\ldots,f_s)\).
Every flat deformation of \(B\) over a local Artin \(k\)-algebra
lifts through any surjection of local Artin \(k\)-algebras.
The tangent space of isomorphism classes of marked deformations
is the cokernel of the Jacobian map
\[
 B^n \longrightarrow B^s,\qquad
 (u_j)_j\longmapsto
 \Big(\sum_j \frac{\partial f_i}{\partial x_j}u_j\Big)_i .
\]
\end{theorem}
\begin{proof}
If \(B_A\) is a flat deformation, choose variable lifts to get
\(P_A=A[x_1,\ldots,x_n]\to B_A\). It is surjective because it is
so modulo \(\mathfrak m_A\) and this ideal is nilpotent. In the
surjection's defining sequence, flatness of \(B_A\) injects
the modulus-\(\mathfrak m_A\) kernel into \(P\).
Thus lifts \(g_i\) of \(f_i\) generate the full kernel:
the quotient by their span is equal to its multiple by
\(\mathfrak m_A\) and hence is zero. For \(A'\twoheadrightarrow A\),
choose arbitrary lifts \(g_i'\in P_{A'}\). Consider their Koszul
complex filtered by powers of \(\mathfrak m_{A'}\).
Every graded quotient complex is a copy of the Koszul complex on
\(f\) with \(k\)-vector-space coefficients \(\mathfrak m_{A'}^j/
\mathfrak m_{A'}^{j+1}\). Its positive homology vanishes by
\cite[RK-07]{core}. Induct along the finite filtration using the
long exact sequence to identify the associated graded of degree-zero
homology with
\(\operatorname{gr}_{\mathfrak m_{A'}}(A')\otimes_k B\).
The filtration on degree zero is by the actual maximal-ideal powers
because it is the image of the filtered degree-zero term of the
Koszul complex. Lift a \(k\)-basis of \(B\) to the quotient.
The resulting map from the free \(A'\)-module is an isomorphism
on associated graded and hence an isomorphism by finite filtration.
This proves flatness, and the presentation gives reduction to \(B_A\).

For first-order deformations use Result~\ref{res:T03} with
\(A=k\), \(A'=k[\epsilon]\), and basepoint the trivial extension.
The classes form the group of Result~\ref{res:T02} for \(N=B\).
For a regular sequence the classes of \(f_i\) freely generate
\(J/J^2\). In a relation \(\sum_i a_i f_i\in J^2\), subtract
a presentation of its right hand side with coefficients in \(J\).
The resulting ordinary relation is a Koszul syzygy, so the remaining
coefficients also lie in \(J\).
The conormal map differentiates the \(f_i\), giving exactly
the displayed Jacobian.
\end{proof}

\begin{example}[A hypersurface tangent space]\label{res:T05}
For \(B=k[x]/(x^m)\), \(m\geq1\), the tangent space is
\(B/(m x^{m-1})\). An infinitesimal deformation is given by
\(x^m+\epsilon q(x)\) and two \(q\)'s give the same class precisely
when their difference vanishes in this quotient.
\end{example}
\begin{proof}
This is the single-entry Jacobian of Result~\ref{res:T04}.
For the polynomial form, the deformed principal equation is
\(x^m+\epsilon q(x)\). Replacing \(x\) by \(x+\epsilon u(x)\)
adds \(m x^{m-1}u(x)\), and multiplying the equation by
\(1+\epsilon v(x)\) adds a multiple of \(x^m\) to \(q\).
These account for the cokernel, including the case when the
characteristic divides \(m\).
\end{proof}

\begin{example}[A first-order deformation that does not lift]\label{res:T06}
Let \(B=k\oplus kx\oplus ky\oplus kz\), with the span of \(x,y,z\)
square-zero. On the free \(k[t]/(t^2)\)-module with this marked basis,
set \(xy=tz\), \(z^2=ty\), and all other products of basis elements
in the maximal-ideal span equal to zero, imposing symmetry and the unit.
This flat first-order deformation admits no flat lift to \(k[t]/(t^3)\).
\end{example}
\begin{proof}
Associativity on triples from the span of \(x,y,z\) holds modulo
\(t^2\), since a product is \(t\) times a member of that span
and the next product adds another factor of \(t\).
If a lift existed, Result~\ref{res:T03} and the free basis argument
allow lifts of the same basis. Up to terms of order three write
\(xy=tz+t^2a\), \(yz=t^2b\), \(z^2=ty+t^2c\).
On the residual algebra, multiplying \(a\) by \(z\) is its
scalar coefficient times \(z\), and multiplying \(b\) by \(x\)
is its scalar coefficient times \(x\). Hence
\[
 (xy)z-x(yz)=t^2\bigl(y+a_0z-b_0x\bigr),
\]
which is nonzero in the claimed free basis. This contradicts
associativity and is independent of the choices of lifted basis.
\end{proof}

\section{Tangent spaces and formal hulls}

Write \(\mathcal A_k\) for the local Artin \(k\)-algebras with residue
field \(k\) and local \(k\)-algebra maps. A small extension has
one-dimensional kernel \(I\) killed by the source maximal ideal.
Every surjection in \(\mathcal A_k\) decomposes into such extensions.
a nonzero kernel has a one-dimensional socle ideal, and induction on
length applies after quotienting it out. Let \(F:\mathcal A_k\to
\mathrm{Sets}\) have \(F(k)=\{*\}\). For \(A_1\to A\leftarrow A_2\)
the canonical fibre-product map is
\[
 F(A_1\times_A A_2)\longrightarrow F(A_1)\times_{F(A)}F(A_2).
\]
We use the classical conditions H1 (surjectivity when \(A_2\to A\)
is small), H2 (bijectivity when \(A=k\), \(A_2=k[\epsilon]\)),
H3 (finite-dimensional tangent space), and H4 (bijectivity for
\(A_1=A_2\) and both maps the same small extension).

\begin{proposition}[Tangent functor and split corrections]\label{res:T07}
Under H2, \(TF=F(k[\epsilon])\) has a natural \(k\)-vector-space structure,
and \(F(k\oplus I)\simeq TF\otimes_k I\) for finite-dimensional
square-zero \(I\). Suppose \(B\to A\) is small and
\(h:A_1\to B\) lifts \(A_1\to A\). Then
\[
 A_1\times_A B \simeq A_1\times_k (k\oplus I)
\]
by \((u,v)\mapsto (u,\bar u+(v-h(u)))\).
Under H1, two prescribed compatible elements of \(F(A_1),F(B)\)
can therefore be represented by an element of \(F(A_1)\) and a
single tangent correction in \(TF\otimes I\).
\end{proposition}
\begin{proof}
Repeated H2 identifies \(F(k\oplus k^r)\) with \(TF^r\).
Apply \(F\) to addition \(k^2\to k\) and multiplication by scalars
on \(k\). The vector-space axioms follow from the corresponding
equalities of linear maps, using the product identifications.
In a basis of \(I\), this is exactly \(TF\otimes I\).
Changing the basis uses the same linear maps under the product
identifications, so the correspondence is natural.
For the splitting formula, the inverse sends
\((u,\bar u+w)\) to \((u,h(u)+w)\). In multiplying, \(w w'=0\)
and the action on \(I\) is via the residue field, which verifies
the formula as a ring isomorphism. Apply H1 to choose an element
over the fibre product and apply H2 to its split presentation.
\end{proof}

\begin{lemma}[Largest quotient on which a lift exists]\label{res:T08}
Suppose H1 holds and \(B\twoheadrightarrow A\) has kernel \(I\)
killed by the maximal ideal of \(B\). For \(x\in F(A)\) there
is a smallest \(k\)-subspace \(J\subseteq I\) such that \(x\)
lifts to \(F(B/J)\).
\end{lemma}
\begin{proof}
The set is nonempty since it contains \(I\) and is upward closed.
By factoring a surjection into small extensions, H1 also gives
surjectivity of the fibre-product map when the second leg is an
arbitrary surjection. For that induction, pull the first ring
back successively along the small extensions. At each stage the
already chosen element and the next prescribed image agree over
the preceding ring, so H1 supplies the next element.
If \(J_1,J_2\) allow lifts, choose \(J_1'\supseteq J_1\) with
\(J_1'+J_2=I\), \(J_1'\cap J_2=J_1\cap J_2\), by adjoining a
vector-space complement to \(J_1+J_2\) in \(I\).
Push the lift forward to \(B/J_1'\).
The two lifts glue to \(B/J_1'\times_A B/J_2=B/(J_1\cap J_2)\).
Thus the set is closed under pairwise intersections.
An element of smallest dimension is contained in every element.
\end{proof}

For a complete Noetherian local \(k\)-algebra \(R\) with residue
field \(k\), put \(h_R(A)=\operatorname{Hom}_{k,\mathrm{loc}}(R,A)\).
A compatible family of elements of \(F(R/\mathfrak m_R^n)\) defines
a natural map \(h_R\to F\). A hull has a tangent isomorphism and the
lifting property that \(h_R(B)\to h_R(A)\times_{F(A)}F(B)\)
is surjective for every small \(B\to A\).

\begin{theorem}[Schlessinger's hull in the classical base category]\label{res:T09}
If H1, H2, and H3 hold, \(F\) has a hull. Conversely the
existence of a hull implies these conditions.
\end{theorem}
\begin{proof}
Let \(S=k[[X_1,\ldots,X_r]]\), \(r=\dim TF\), with maximal ideal
\(\mathfrak n\). Set \(J_2=\mathfrak n^2\). The element
\(\xi_2\in F(S/J_2)\) corresponding under Result~\ref{res:T07}
to the tensor of a basis of \(TF\) with the variable classes induces
a tangent isomorphism. Having chosen \(\xi_n\in F(S/J_n)\), apply
Result~\ref{res:T08} to
\(S/\mathfrak n J_n\to S/J_n\) and choose the smallest ideal
\(J_{n+1}\) between \(\mathfrak n J_n\) and \(J_n\) through which
\(\xi_n\) lifts. Choose a lift \(\xi_{n+1}\).
All these quotients are Artin since \(\mathfrak n^n\subseteq J_n\)
by induction.

Put \(J=\bigcap J_n\) and \(R=S/J\). The ideals \(J_n/J\) are
cofinal with powers of the maximal ideal of \(R\).
For the nontrivial direction, for each \(q\) the images of \(J_n\) in
\(S/\mathfrak n^q\) stabilize. Their stable images surject to one
another as \(q\) decreases by choosing \(n\) large in two successive
degrees. A point in a stable image therefore lifts to compatible
stable points in all higher degrees, hence to an element of \(S\).
For each fixed \(n\) this element lies in \(J_n+\mathfrak n^q\)
for every \(q\), and thus in \(J_n\) by Krull intersection on
\(S/J_n\). It lies in \(J\). The stable image is precisely the
image of \(J\), as required. The compatible \(\xi_n\) therefore
give a natural map \(h_R\to F\) with tangent isomorphism.

To prove the lifting property let \(B\to A\) be small with kernel
\(I\), \(a:R\to A\), and \(y\in F(B)\) lifting the element from \(a\).
Factor \(a\) through \(S/J_n\). Consider \(T=S/J_n\times_A B\).
If the projection \(T\to S/J_n\) splits, composing its section
with \(R\to S/J_n\) gives a ring lift \(R\to B\).
If it does not split, any map \(S\to T\) lifting the quotient map
\(S\to S/J_n\) is surjective. Indeed its image surjects onto
\(S/J_n\), and if it missed the one-dimensional kernel it would
be a section. Such a map exists by lifting the \(X_i\) to the
maximal ideal of \(T\). Its kernel \(K\) lies between
\(\mathfrak n J_n\) and \(J_n\), since the kernel of the projection
is killed by the maximal ideal. H1 glues \(\xi_n,y\) to an
element over \(T=S/K\). Thus \(J_{n+1}\subseteq K\) by
Result~\ref{res:T08}, giving again a ring lift \(R\to B\).

It remains to make the ring lift send the formal object to the
prescribed \(y\). Increase \(n\) so that this lift factors through
\(S/J_n\). By H1 glue \(\xi_n,y\) over \(T=S/J_n\times_A B\).
The chosen lift splits the projection, and Result~\ref{res:T07}
expresses the glued element by \((\xi_n,\theta)\) for a tangent
correction \(\theta\in F(k\oplus I)\). The tangent isomorphism
gives \(\delta:R\to k\oplus I\) carrying the formal object to
\(\theta\). The map \(R\to (S/J_n)\times_k(k\oplus I)\) with
components the quotient and \(\delta\) sends it to the glued
element by H2. Its composition through \(T\to B\) gives the
required lift matching \(y\).

For the converse, smoothness of a hull and \(F(k)=\{*\}\) make
\(h_R(A)\to F(A)\) surjective by induction on length.
To glue compatible elements for a small leg, choose a ring map
for the first element and use smoothness to choose one for the
second with the same reduction. Their fibre-product map proves H1.
For H2, the same existence argument gives surjectivity. For
injectivity, take any two elements over \(A_1\times_k k[\epsilon]\)
with the same projections, and lift them from \(h_R\).
Their tangent components agree since the tangent map is injective.
One may instead fix a single ring lift of the \(A_1\)-component.
Smoothness for the small projection
\(A_1\times_k k[\epsilon]\to A_1\) then lifts both elements with
that fixed component. Now the ring maps have both components
equal, so the elements coincide. H3 follows from Noetherianity
of \(R\) and the tangent isomorphism.
\end{proof}

\begin{theorem}[Prorepresentability]\label{res:T10}
Under H1--H3, the hull represents the functor on \(\mathcal A_k\)
if and only if H4 holds. Conversely every prorepresentable
functor in this category satisfies H1--H4.
\end{theorem}
\begin{proof}
Let \(B\to A\) be small and fix \(x\in F(A)\). When its fibre
of lifts is nonempty, H4 identifies ordered pairs of lifts with
the elements of \(F(B\times_A B)\) over \((x,x)\).
Use the split formula of Result~\ref{res:T07} with \(h=1_B\).
With a first lift fixed, the second is therefore in bijection
with \(TF\otimes I\). For three lifts, apply H1 to
\((B\times_A B)\to A\leftarrow B\) to glue the first pair
and the third lift over the triple fibre product. Its pairwise
restrictions are the H4 representatives of the three pairs.
Their two successive differences add to the first-to-last
difference. The ring maps
from the triple fibre product to \(k\oplus I\) send
\((v_1,v_2,v_3)\) to the residue with components \(v_2-v_1\),
\(v_3-v_2\), and \(v_3-v_1\), respectively. H2 and the addition
rule of Result~\ref{res:T07} transport that equality to the
tangent space. Hence the fibre is a torsor, naturally in \(F\).

If \(h_R\to F\) is a hull, it is surjective over each small step.
Inductively assume it bijects elements over the quotient \(A\).
The nonempty lift fibres over corresponding elements are torsors
on both sides, using H4 for \(F\) and the fibre-product property
of \(h_R\). The tangent isomorphism makes their equivariant
maps bijective. Induction on length gives the natural isomorphism
\(h_R\simeq F\). Conversely \(\operatorname{Hom}(R,-)\) preserves
fibre products, and its tangent space is the dual of the
finite-dimensional cotangent space of \(R\), giving all conditions.
\end{proof}

The obstruction in Result~\ref{res:T01} is a concrete extension
test on a conormal module, and Results~\ref{res:T02} and~\ref{res:T03}
give its ambiguity and flatness tests. For general cotangent
complexes, the derived obstruction can be packaged as degree-two
Ext, with ambiguity in degree one and infinitesimal automorphisms
in degree zero. The full construction and its comparison are
treated in \cite[\S16]{st-cot}.
